% p.59 — الباب الخامس والعشرون في التصريف مع الضمائر وفي اللازم والمتعدي
% Layout (as in the lithograph): one `tables` row. RIGHT: five narrow columns of 21 cells
% each (right to left), read top to bottom, column by column; the lithograph rules off groups
% of related forms with horizontal lines. LEFT: a sixth, wider column of prose (definitions of
% اللازم and المتعدي and the rule on making a lāzim verb mutaʿaddī), set as \sidetext beside
% the table.
% Initial hamza is printed as bare alif + vowel (اِنَّهُ، اِيَّاهُ); transcribed with hamza.
\mspage{59}
\bab{الباب الخامس والعشرون في التصريف مع الضمائر وفي اللازم والمتعدي}
\begin{tables}
\begin{sarf}{|c|c|c|c|c|}
إِنَّهُ & إِيَّاكُمَا & عَلَيْنَا & لَيْتَنِيْ & فِيْهِ\\
إِنَّهُمَا & إِيَّاكُمْ & إِلٰى & لَيْتَنَا & فِيَّ\\
إِنَّهُمْ & إِيَّاكِ & إِلَيْهِ & لَكِنَّهُ & رُبَّهُ\\
إِنَّهَا & إِيَّاكُمَا & إِلَيَّ & لَكِنِّيْ & رُبَّنِيْ\\
إِنَّهُمَا & إِيَّاكُنَّ & إِلَيْنَا & لَكِنَّا & رُبَّنَا\\
إِنَّهُنَّ & إِيَّايَ & فِيْ & بَيْنَهُ & دَارِيْ\\
إِنَّكَ & إِيَّانَا & فِيْهِ & بَيْنِيْ & ثَوْبِيْ\\
إِنَّكُمَا & عَلٰى & فِيَّ & بَيْنَنَا & كِتَابِيْ\\
إِنَّكُمْ & عَلَيْهِ & فِيْنَا & لَدٰى & غُلَامِيْ\\
إِنَّكِ & عَلَيْهِمَا & بِ & لَدَيْهِ & أُمِّيْ\\
إِنَّكُمَا & عَلَيْهِمْ & بِهِ & لَدَيَّ & قَلَمُكَ\\
إِنَّكُنَّ & عَلَيْهَا & بِيْ & لَدَيْنَا & فَوْقَهُ\\
إِنَّنِيْ & عَلَيْهِمَا & بِنَا & أَخُوْهُ & تَحْتَهُ\\
إِنَّنَا & عَلَيْهِنَّ & لِ & أَخَاهُ & يَمِيْنَهُ\\
إِيَّاهُ & عَلَيْكَ & لَهُ & أَخِيْهِ & يَسَارَهُ\\
إِيَّاهُمَا & عَلَيْكُمَا & لِيْ & أَخِيْ & أَمَامَهُ\\
إِيَّاهُمْ & عَلَيْكُمْ & لَنَا & أَبُوْهُ & وَرَاءَهُ\\
إِيَّاهَا & عَلَيْكِ & لَعَلَّهُ & أَبَاهُ & قَبْلَهُ\\
إِيَّاهُمَا & عَلَيْكُمَا & لَعَلِّيْ & أَبِيْهِ & بَعْدَهُ\\
إِيَّاهُنَّ & عَلَيْكُنَّ & لَعَلَّنَا & فُوْهُ & قَبْلَ زَيْدٍ\\
إِيَّاكَ & عَلَيَّ & لَيْتَهُ & فَاهُ & دَارَ زَيْدٍ\\
\end{sarf}
% Left-hand prose column:
\sidetext{0.3\linewidth}{%
اللازم = فعل لا يطلب المفعول، ذَهَبَ زَيْدٌ.
\par
المتعدى = فعل يطلب المفعول، ضَرَبَ زَيْدٌ عَمْرًا.
\par
لا يؤخذ من الفعل اللازم الماضي المجهول والمضارع المجهول والمصدر المجهول واسم المفعول، مثل ذَهَبَ يَذْهَبُ ذَهَابًا ذَاهِبٌ اِذْهَبْ لَا تَذْهَبْ.
\par
و اذا اردت ان تجعل اللازم متعديا فتأتى بالهمزة مثل خَرَجَ أَخْرَجَ، او بالتشديد مثل بَلَغَ بَلَّغَ، او بحرف الجر مثل أَتٰى أَتٰى بِهِ.
}
\end{tables}
